%%%%%%%%%%%%%%%%%%%%%%%%%%%%%%%%%%%%%%%%%%%%%%%%%%%%%%%%%%%%%%%%%%%%%%%%%%%
% Lehreinheit 9: Qualitative Methoden
% Quelle: WuZ-Buch, Kap. 15 (LE 9), Kap. 6 (Persona, CREDS,
%         Re-Enchantment), Kap. 5b (Vogelmann: Wahrheitspraktiken),
%         Abschnitt 10f (Ideomotor-Pendel), Kap. 10 (Charpentier)
% Staunmoment: Das Patter -- derselbe Effekt zweimal, sachlich und
%              magisch erzählt
% Kernbegriff: Intersubjektivität
%%%%%%%%%%%%%%%%%%%%%%%%%%%%%%%%%%%%%%%%%%%%%%%%%%%%%%%%%%%%%%%%%%%%%%%%%%%
\nsl{LE~8 hat gezählt: Gewinne beim Wechseln, Ausschläge in Millimetern,
Staunen auf einer Skala von~1 bis~5. Diese Einheit fragt, was dabei
verloren geht. Sie beginnt mit einem Effekt, den die Studierenden schon
kennen, und zeigt ihn zweimal mit identischen Handlungen. Nur die Sprache
ändert sich: einmal sachlich, einmal mit dem \enquote{magischen}
Sprechtext, den Zauberkünstler \emph{Patter} nennen. Was sich dabei ändert,
lässt sich zwar auf einer Skala ankreuzen, aber erst in den Worten der
Beteiligten verstehen. Daraus entwickelt die Einheit die Grundlogik
qualitativer Forschung: Verstehen statt Messen, Wirklichkeit als
sprachlich mitkonstruiert, Objektivität als Intersubjektivität. Sie
schließt mit zwei Werkzeugen, die die Studierenden sofort anwenden: dem
Leitfadeninterview und der qualitativen Inhaltsanalyse.
\newpage}
\os{\begin{description}
\item[{\rot\bf Heute}]~\\[-\lblskip]
\bi
  \ite derselbe Effekt zweimal -- nur die Worte ändern sich
  \ite Verstehen statt Messen
  \ite Konstruktivismus und Intersubjektivität
  \ite Interview, Fokusgruppe, Beobachtung, Dokumente
  \ite Gütekriterien und qualitative Inhaltsanalyse
\ei
\end{description}
\newpage}

\begin{vorbereitung}[Das Patter: ein Effekt, zwei Sprechtexte]
\textbf{Eine Woche vorher}
\begin{checkliste}
  \item \textbf{Effekt wählen.} Zwei Varianten sind vorbereitet, beide
    aus dem Buch und beide im Semester schon vorgekommen:
    das \emph{Ideomotor-Pendel} (LE~1; Buch, Abschnitt~10f) oder die
    \emph{Charpentier-Illusion} (LE~7; Buch, Kap.~10). Empfohlen ist das
    Pendel, weil alle gleichzeitig mitmachen und der Effekt stark auf
    Erwartung reagiert. Die Charpentier-Illusion eignet sich, wenn der
    Raum kein Aufstützen erlaubt; dann heben 6--8 Freiwillige vorn,
    der Rest beobachtet und notiert.
  \item Dass die Studierenden den Mechanismus kennen, ist Absicht: Das
    Wissen (Belief) bleibt in beiden Runden gleich, nur die Sprache
    ändert sich. Verändert sich das Erleben trotzdem, wirkt die Sprache
    auf der Ebene der Alief (LE~1).
  \item \textbf{Material Pendel:} die Pendel aus LE~1 (Faden 40\,cm,
    Mutter M8 oder Ring), zehn Ersatzpendel.
    \textbf{Material Charpentier:} drei gleich aussehende
    Streichholzschachteln mit Gummiband, die oberste mit ca.\ 50\,g
    beschwert (wie in LE~7); zwei Sets.
  \item \textbf{Beide Sprechtexte} (Abschnitt Staunmoment) so oft laut
    üben, bis sie ohne Ablesen sitzen. Der magische Text lebt von
    Pausen, leiser Stimme und Blickkontakt, der sachliche von gleichmäßigem
    Tempo und Laborton. Die Handlungen müssen in beiden Runden identisch
    sein: gleiche Haltung, gleiche Dauer (Stoppuhr, je 60 Sekunden
    Pendeln), gleiche Kommandos für die Richtung.
  \item Pro Person zwei Karteikarten (A7) vorbereiten, beschriftet mit
    \enquote{Runde~1} und \enquote{Runde~2}. Darauf kommen nach jeder
    Runde eine Zahl (Stärke des Erlebens, 1--5) und drei Wörter.
\end{checkliste}

\textbf{Handout und Übungsblatt}
\begin{checkliste}
  \item Handout \enquote{Qualitative Methoden} und Übungsblatt LE~9 je
    einmal pro Person kopieren (Druckvorlage \texttt{handouts.pdf},
    Seiten zu LE~9). Das Übungsblatt enthält die Leitfaden-Vorlage, das
    konstruierte Interviewtranskript und die Codiertabelle.
  \item Für den Codier-Vergleich das Transkript zusätzlich einmal pro
    Gruppe auf A3 kopieren, damit Gruppen Segmente markieren können;
    zwei Farben Textmarker pro Gruppe.
\end{checkliste}

\textbf{Am Tag selbst}
\begin{checkliste}
  \item Tafel teilen: links \enquote{Runde~1}, rechts \enquote{Runde~2};
    darunter je ein Feld für die Zahlen und eines für die Wörter.
  \item Licht für Runde~2 etwas dimmen, falls möglich; das ist Teil des
    Patters und wird im Debrief offengelegt.
  \item Karten nach jeder Runde sofort einsammeln, damit die Antworten
    der zweiten Runde nicht an die der ersten angepasst werden.
\end{checkliste}

\textbf{Risiken und Plan B.} Manche Studierende reagieren auf den
magischen Text mit Ironie oder Abwehr. Das ist kein Störfall, sondern
Material: Auch Abwehr ist eine Reaktion auf Sprache und gehört auf die
Karte. Die Reihenfolge (erst sachlich, dann magisch) ist mit Übung und
Ermüdung vermengt; das wird im Debrief als Designproblem besprochen
(Brücke zu LE~6). Ethisch gilt, was das Buch für die Persona der
Lehrenden festhält: keine falschen Überzeugungen pflanzen. Der magische
Rahmen wird deshalb noch in derselben Viertelstunde als Rahmen
offengelegt. Fällt der Effekt ganz aus (Pendel schwingt kaum), wird
direkt mit den Karten gearbeitet: Auch \enquote{nichts gespürt} ist eine
Beschreibung, deren Wortwahl sich zwischen den Runden unterscheiden kann.
\end{vorbereitung}

\begin{ablauf}
\begin{description}
\item[0--5 · Blitzlicht.] Anknüpfen an LE~8: Nicht alles Wichtige ist
  zählbar. Was messen Zahlen \emph{nicht}?
\item[5--20 · Staunmoment.] Derselbe Effekt zweimal, erst mit sachlichem,
  dann mit magischem Sprechtext. Nach jeder Runde Karte: eine Zahl,
  drei Wörter. Zahlen und Wörter an die Tafel; Offenlegung des Rahmens.
\item[20--50 · Theorie-Input.] Verstehen statt Messen; Konstruktivismus
  und Intersubjektivität; Patter und Wahrheitspraktiken (Buch, Kap.~6
  und 5b); Methodenüberblick; Leitfadeninterview; Gütekriterien;
  qualitative Inhaltsanalyse nach Mayring.
\item[50--75 · Interview-Design.] Leitfaden mit 8--12 Fragen zur eigenen
  Forschungsfrage (Übungsblatt, Aufgabe~2); Probe-Interview in Paaren;
  Reflexion suggestiver und offener Fragen.
\item[75--90 · Codier-Vergleich.] Zwei Gruppen codieren dasselbe kurze
  Transkript (Aufgabe~3) und vergleichen ihre Codes.
\end{description}
\end{ablauf}

%==============================================================================
\section{Staunmoment: Zweimal dasselbe}
%==============================================================================
\begin{staunbox}{Das Patter}
Sie erleben denselben Effekt zweimal. Die Handlungen sind identisch:
gleiche Haltung, gleiche Dauer, gleiche Richtungen.

\emph{Nach jeder Runde notieren Sie auf Ihrer Karte: eine Zahl von 1 bis
5 (wie stark war Ihr Erleben?) und drei Wörter, die es beschreiben.}
\end{staunbox}
\os{\newpage}

\begin{lehrhinweis}[Sprechtext A: sachlich (Pendel, Runde 1)]
\itshape
Nehmen Sie das Pendel zur Hand. Ellbogen auf den Tisch, Faden zwischen
Daumen und Zeigefinger, das Gewicht etwa zwei Zentimeter über der
Handfläche. Wir machen eine Messung. Ich gebe eine Richtung vor, Sie
stellen sich die Bewegung nur vor. Richtung eins: links -- rechts. Ich
stoppe nach sechzig Sekunden. \dots\ Danke. Richtung zwei: vor --
zurück. \dots\ Danke. Legen Sie das Pendel ab und notieren Sie auf Karte
eins eine Zahl und drei Wörter.
\end{lehrhinweis}

\begin{lehrhinweis}[Sprechtext B: magisch (Pendel, Runde 2)]
\itshape
Nehmen Sie das Pendel noch einmal. Diesmal halten Sie es anders -- nicht
mit den Fingern, sondern mit Aufmerksamkeit. Seit Jahrhunderten haben
Menschen Pendeln Fragen gestellt, deren Antwort sie nicht aussprechen
wollten. Ihr Kopf weiß manches erst später als Ihre Hand. \dots\
Ellbogen auf den Tisch. Lassen Sie das Gewicht zur Ruhe kommen. Und jetzt
fragen Sie es, ohne die Lippen zu bewegen: Wohin willst du? Wenn es nach
links und rechts geht \dots\ lassen Sie es gehen. \dots\ Und jetzt bitten
Sie es: vor -- und zurück. Merken Sie, wie es zögert, bevor es Ihnen
folgt? \dots\ Legen Sie es behutsam ab. Karte zwei: eine Zahl, drei Wörter.
\end{lehrhinweis}

\begin{lehrhinweis}[Variante Charpentier-Illusion]
\textbf{Sachlich:} \emph{Drei Schachteln, mit Gummiband verbunden. Heben
Sie den Stapel an, stellen Sie ihn ab. Nehmen Sie nur die oberste
Schachtel und heben Sie sie an. Schätzen Sie: leichter, gleich oder
schwerer als der Stapel?} \textbf{Magisch:} \emph{Diese drei Schachteln
habe ich aus dem Nachlass eines Uhrmachers. Er hat in die oberste etwas
gelegt, das nie jemand gesehen hat. Man sagt, sie werde schwerer, sobald
sie allein ist -- als vermisse sie die anderen. Heben Sie den Stapel. Und
jetzt nehmen Sie sie aus der Gemeinschaft heraus \dots} Der magische Text
nutzt die drei Mechanismen, mit denen Objekte nach dem Buch
\enquote{geladen} werden: Geschichte, Zeremonie der Einführung, Teilhabe
des Publikums (Buch, Abschnitt Re-Enchantment; \textcite{taylor2020bizarre}).
Die Geschichte ist erfunden und wird im Debrief als erfunden benannt.
\end{lehrhinweis}

\begin{description}
\item[{\rot\bf Was war anders?}]~\\[-\lblskip]
\bi
  \ite Zahlen an die Tafel: Runde~1 und Runde~2
  \ite Wörter an die Tafel: Runde~1 und Runde~2
  \ite die Handlungen waren identisch -- was hat sich geändert?
  \ite was sagen Ihnen die Zahlen, was sagen Ihnen die Wörter?
\ei
\nsl{Die Zahlen zeigen bestenfalls, \emph{dass} sich etwas geändert hat
und in welche Richtung. Die Wörter zeigen, \emph{wie}: Denkbar sind
nach der sachlichen Runde Wörter wie \enquote{Messung} oder
\enquote{gleichmäßig}, nach der magischen eher \enquote{unheimlich},
\enquote{albern} oder \enquote{als ob jemand zieht}. Was die eigene
Gruppe tatsächlich schreibt, zeigt erst die Tafel; Vorhersagen dazu
sollte man im Raum nicht äußern. Wichtig ist der Hinweis, dass auch
\enquote{albern} eine Wirkung der Sprache ist.}
\os{\newpage}

\item[{\rot\bf Offenlegung}]~\\[-\lblskip]
\bi
  \ite das Pendel hat in beiden Runden dasselbe getan: Ihre Hand bewegt
    es (LE~1)
  \ite geändert haben sich Wortwahl, Tempo, Licht, Stimme -- der
    \emph{Rahmen}
  \ite und: die Reihenfolge -- Runde~2 war immer die geübtere
  \ite Frage für heute: Wie erforscht man, was ein Rahmen mit Menschen
    macht?
\ei
\os{\newpage}
\end{description}

%==============================================================================
\section{Verstehen statt Messen}
%==============================================================================
\nsl{Quantitative Forschung fragt, wie viel und wie oft; sie zählt
Merkmale, die vorab festgelegt und operationalisiert wurden (LE~2).
Qualitative Forschung fragt, wie und warum; sie will verstehen, welche
Bedeutung Menschen einer Situation geben, und entwickelt ihre Kategorien
oft erst am Material. Die beiden Zugänge sind keine Konkurrenten. Die
Zahlen auf den Karten haben gezeigt, dass sich das Erleben verschoben
hat; die Wörter haben gezeigt, worin. Für die zweite Frage braucht man
offene Daten und Verfahren, die mit Sprache umgehen können.}
\begin{description}
\item[{\rot\bf Zwei Zugänge}]~\\[-\lblskip]
\begin{center}
\small
\begin{tabular}{@{}p{0.2\textwidth}p{0.33\textwidth}p{0.33\textwidth}@{}}
\toprule
 & \textbf{quantitativ} & \textbf{qualitativ} \\
\midrule
Frage & wie viel, wie oft, ob & wie, warum, was bedeutet es \\
Daten & Zahlen, standardisiert & Texte, Bilder, Beobachtungen \\
Fälle & viele & wenige, gezielt ausgewählt \\
Kategorien & vorab festgelegt & oft am Material entwickelt \\
Ziel & prüfen, verallgemeinern & verstehen, Hypothesen bilden \\
\bottomrule
\end{tabular}
\end{center}
\os{\newpage}

\item[{\rot\bf Was Zahlen nicht messen}]~\\[-\lblskip]
\bi
  \ite \enquote{4} auf der Karte kann heißen: beeindruckt -- oder
    beunruhigt
  \ite eine Skala erfasst nur, was jemand vorher gefragt hat
  \ite Bedeutungen, Begründungen, Widersprüche fallen durch das Raster
  \ite qualitative Daten zeigen, \emph{welche} Fragen man später zählen
    sollte
\ei
\os{\newpage}
\end{description}

%==============================================================================
\section{Konstruktivismus und Intersubjektivität}
%==============================================================================
\nsl{Der Patter-Versuch macht eine These erfahrbar, die die
Sozialwissenschaften unter dem Stichwort \emph{Konstruktivismus}
diskutieren. Nach \textcite{bergerluckmann1966} bringen Menschen ihre
gemeinsame Alltagswirklichkeit in Interaktion und Sprache hervor; was so
entstanden ist, erscheint ihnen anschließend als objektiv gegeben. Die
Zauberkunst arbeitet bewusst mit diesem Mechanismus. Das Buch beschreibt,
wie Kontext und Persona darüber entscheiden, ob ein Effekt als Trick oder
als Wunder erlebt wird (Kap.~6), und wie Objekte durch Geschichte,
Zeremonie und Teilhabe \enquote{geladen} werden
\parencite{taylor2020bizarre}. Frieder Vogelmann fragt in diesem Sinn
nicht nur, ob eine Aussage wahr ist, sondern wie sie \emph{wirksam}
geworden ist: durch welche Wahrheitspraktiken, Inszenierungen und
Autoritäten \parencite{vogelmann2020wirksamkeit}. Das Buch folgert daraus,
dass Wissenschaft sich von der inszenierten Wahrheit des Zauberkünstlers
durch Transparenz und Intersubjektivität abheben muss: Sie legt ihr
Verfahren offen, wo der Zauberkünstler es verbirgt.}
\begin{description}
\item[{\rot\bf Sprache konstruiert Erleben}]~\\[-\lblskip]
\bi
  \ite das Patter verändert nicht das Pendel, sondern das, was Sie
    erleben
  \ite Kontext und Persona entscheiden, ob ein Effekt Trick oder Wunder
    ist (Buch, Kap.~6)
  \ite Wahrheit wird \emph{wirksam} durch Praktiken, Rollen, Inszenierung
    \parencite{vogelmann2020wirksamkeit}
  \ite Wissenschaft hält dagegen: Verfahren offenlegen statt verbergen
\ei
\os{\newpage}

\item[{\rot\bf Jedes Interview ist auch ein Patter}]~\\[-\lblskip]
\bi
  \ite Unfallfilm, Frage nach der Geschwindigkeit: \enquote{smashed}
    statt \enquote{hit} führt zu höheren Schätzungen
  \ite eine Woche später erinnern mehr Personen zerbrochenes Glas, das es
    nicht gab \parencite{loftuspalmer1974}
  \ite Fragen erheben Wirklichkeit nicht nur, sie formen sie mit
  \ite Konsequenz: die eigene Sprache als Teil der Daten behandeln
\ei
\os{\newpage}

\item[{\rot\bf Intersubjektivität}]~\\
\defn{intersubjektivitaet}{\textbf{Intersubjektivität} heißt: Ein Befund
hängt nicht an der Person, die ihn erhoben hat. Andere können das
Vorgehen nachvollziehen und kommen mit demselben Material zu
vergleichbaren Ergebnissen oder können begründen, warum nicht.}
\bi
  \ite in der quantitativen Forschung: standardisierte Messung
  \ite in der qualitativen Forschung: offengelegte Regeln, mehrere
    Codierende, dokumentierte Entscheidungen
\ei
\os{\newpage}
\end{description}

%==============================================================================
\section{Methoden im Überblick}
%==============================================================================
\nsl{Qualitative Daten entstehen auf vier Hauptwegen. Im
\emph{Leitfadeninterview} spricht eine Person entlang vorbereiteter, aber
offener Fragen; der Leitfaden sichert Vergleichbarkeit, ohne die
Antworten vorzugeben. In der \emph{Fokusgruppe} diskutieren sechs bis
zehn Personen ein Thema; sichtbar werden dabei geteilte Deutungen und
Konflikte, die im Einzelgespräch nicht entstünden. Die
\emph{Beobachtung} (teilnehmend oder nicht teilnehmend, offen oder
verdeckt) erfasst, was Menschen tun, nicht nur, was sie darüber sagen.
Die \emph{Dokumentenanalyse} arbeitet mit vorhandenem Material wie
Protokollen, Zeitungsartikeln, Lehrplänen oder Social-Media-Beiträgen.
Eine ausführliche Darstellung aller vier Wege bietet
\textcite{doeringbortz2016}.}
\begin{description}
\item[{\rot\bf Vier Wege zu qualitativen Daten}]~\\[-\lblskip]
\begin{center}
\small
\begin{tabular}{@{}p{0.22\textwidth}p{0.33\textwidth}p{0.33\textwidth}@{}}
\toprule
\textbf{Methode} & \textbf{stark bei \dots} & \textbf{Vorsicht bei \dots} \\
\midrule
Leitfadeninterview & Deutungen, Begründungen & Suggestion, sozialer Erwünschtheit \\
Fokusgruppe & geteilten Sichtweisen, Konflikten & Meinungsführern, Gruppendruck \\
Beobachtung & tatsächlichem Verhalten & Beobachtereffekt, Deutung \\
Dokumentenanalyse & Prozessen, Diskursen & Auswahl der Dokumente \\
\bottomrule
\end{tabular}
\end{center}
\os{\newpage}
\end{description}

%==============================================================================
\section{Das Leitfadeninterview}
%==============================================================================
\nsl{Ein Leitfaden ist kein Fragebogen. Er strukturiert das Gespräch in
wenige Themenblöcke, beginnt jeden Block mit einer offenen,
erzählgenerierenden Frage und hält Nachfragen bereit, die erst gestellt
werden, wenn ein Aspekt nicht von selbst angesprochen wird.
\textcite{helfferich2011qualitaet} empfiehlt, Leitfäden in vier Schritten
zu bauen: zunächst alle Fragen sammeln, dann prüfen (Ist die Frage offen?
Will ich das wirklich wissen? Setzt sie Wissen voraus?), sortieren und
schließlich unter wenige erzählgenerierende Leitfragen subsumieren. Das
Interview ist dabei immer eine gemeinsam hergestellte
Kommunikationssituation; Neutralität ist nicht erreichbar, wohl aber
Reflexion über die eigene Rolle.}
\begin{description}
\item[{\rot\bf Einen Leitfaden bauen}]~\\[-\lblskip]
\bi
  \ite \textbf{sammeln}: alle Fragen notieren, ohne zu filtern
  \ite \textbf{prüfen}: offen? nötig? suggestiv? Fachwort?
  \ite \textbf{sortieren}: nach Themenblöcken
  \ite \textbf{subsumieren}: je Block eine offene Leitfrage, darunter
    Nachfragen \parencite{helfferich2011qualitaet}
  \ite 8--12 Fragen reichen für 30 Minuten
\ei
\os{\newpage}

\item[{\rot\bf Offen oder suggestiv?}]~\\[-\lblskip]
\begin{center}
\small
\begin{tabular}{@{}p{0.43\textwidth}p{0.43\textwidth}@{}}
\toprule
\textbf{suggestiv / geschlossen} & \textbf{offen} \\
\midrule
\enquote{Fanden Sie den Trick auch so verblüffend?} & \enquote{Wie haben Sie den Moment erlebt?} \\
\enquote{Hat Ihnen das beim Lernen geholfen?} & \enquote{Was ist Ihnen davon im Kopf geblieben?} \\
\enquote{Warum mögen Sie Statistik nicht?} & \enquote{Erzählen Sie, wie es Ihnen mit Statistik geht.} \\
\bottomrule
\end{tabular}
\end{center}
\bi
  \ite Faustregel: Die Frage darf die Antwort nicht enthalten
\ei
\os{\newpage}
\end{description}

%==============================================================================
\section{Gütekriterien}
%==============================================================================
\nsl{Die klassischen Gütekriterien Objektivität, Reliabilität und
Validität setzen standardisierte Messung voraus. Für interpretative
Forschung haben \textcite{lincolnguba1985naturalistic} vier Entsprechungen
vorgeschlagen, die bis heute verbreitet sind. \emph{Credibility} fragt, ob
die Rekonstruktion für die Beteiligten glaubwürdig ist; Mittel sind
längeres Feldengagement, Triangulation und die Rückmeldung an die
Befragten (\emph{member checking}). \emph{Transferability} fragt, ob
Befunde auf andere Kontexte übertragbar sind; dazu muss der Kontext so
genau beschrieben sein, dass Lesende selbst urteilen können.
\emph{Dependability} fragt nach der Nachvollziehbarkeit des Vorgehens,
\emph{Confirmability} danach, ob die Befunde aus den Daten folgen und
nicht aus den Vorlieben der Forschenden; beides sichert ein
dokumentierter Prüfpfad (\emph{audit trail}).}
\begin{description}
\item[{\rot\bf Gütekriterien nach Lincoln und Guba}]~\\[-\lblskip]
\begin{center}
\small
\begin{tabular}{@{}p{0.2\textwidth}p{0.2\textwidth}p{0.46\textwidth}@{}}
\toprule
\textbf{Kriterium} & \textbf{entspricht} & \textbf{gesichert durch} \\
\midrule
Credibility & interner Validität & Triangulation, member checking \\
Transferability & externer Validität & dichte Beschreibung des Kontexts \\
Dependability & Reliabilität & dokumentiertes Vorgehen \\
Confirmability & Objektivität & Audit Trail, Reflexion der eigenen Rolle \\
\bottomrule
\end{tabular}
\end{center}
\nsl{Quelle: \textcite{lincolnguba1985naturalistic}.}
\os{\newpage}
\end{description}

%==============================================================================
\section{Qualitative Inhaltsanalyse nach Mayring}
%==============================================================================
\nsl{Die qualitative Inhaltsanalyse ist im deutschsprachigen Raum das
verbreitetste Verfahren, um Texte regelgeleitet auszuwerten
\parencite{mayring2015inhaltsanalyse}. Sie zerlegt das Material in
Analyseeinheiten und ordnet diese Kategorien zu. Die Kategorien können
vorab aus der Theorie abgeleitet werden (deduktiv) oder am Material
entstehen (induktiv); in der Praxis wird oft beides kombiniert. Jede
Kategorie erhält eine Definition, ein Ankerbeispiel aus dem Material und
eine Kodierregel für Abgrenzungsfälle. Die Regelgeleitetheit unterscheidet
die Inhaltsanalyse vom bloßen Interpretieren: Eine zweite Person soll mit
demselben Kategoriensystem zu vergleichbaren Zuordnungen kommen. Wie gut
das gelingt, misst die Intercoder-Übereinstimmung.}
\begin{description}
\item[{\rot\bf Ablauf in fünf Schritten}]~\\[-\lblskip]
\begin{center}
\resizebox{0.95\textwidth}{!}{%
\begin{tikzpicture}[font=\sffamily,>=Stealth,
  st/.style={draw=rwuviolet,thick,fill=rwuviolet!8,rounded corners=2pt,
    minimum width=27mm,minimum height=13mm,align=center}]
  \node[st] (a) at (0,0) {Material\\festlegen};
  \node[st] (b) at (3.8,0) {Einheiten\\bestimmen};
  \node[st] (c) at (7.6,0) {Kategorien\\definieren};
  \node[st,draw=rwucyan,fill=rwucyan40!40] (d) at (11.4,0) {codieren};
  \node[st] (e) at (15.2,0) {prüfen und\\auswerten};
  \draw[->,thick] (a) -- (b);
  \draw[->,thick] (b) -- (c);
  \draw[->,thick] (c) -- (d);
  \draw[->,thick] (d) -- (e);
  \draw[->,thick,rwugray,dashed] (d.south) to[bend left=25]
    node[below]{Probedurchlauf: Kategorien überarbeiten} (c.south);
\end{tikzpicture}}
\end{center}
\bi
  \ite deduktiv: Kategorien aus Theorie oder Forschungsfrage
  \ite induktiv: Kategorien aus dem Material
\ei
\os{\newpage}

\item[{\rot\bf Eine Kategorie definieren}]~\\[-\lblskip]
\begin{center}
\small
\begin{tabular}{@{}p{0.2\textwidth}p{0.66\textwidth}@{}}
\toprule
\textbf{Kategorie} & K1 \enquote{Staunen} \\
\midrule
\textbf{Definition} & Äußerungen, in denen die Person beschreibt, dass
etwas ihren Erwartungen widersprach und sie überrascht oder verblüfft war \\
\textbf{Ankerbeispiel} & \enquote{Ich hab wirklich gedacht, das geht
nicht.} \\
\textbf{Kodierregel} & nur codieren, wenn die Überraschung erlebt wird;
bloße Bewertung (\enquote{war gut gemacht}) gehört zu K5 \\
\bottomrule
\end{tabular}
\end{center}
\nsl{Das Beispiel stammt aus dem Kategoriensystem des Übungsblatts.}
\os{\newpage}

\item[{\rot\bf Codieren und vergleichen}]~\\[-\lblskip]
\bi
  \ite zwei Personen codieren dasselbe Material unabhängig
  \ite einfache Übereinstimmung: übereinstimmende Segmente durch alle
    Segmente
  \ite Abweichungen sind kein Fehler, sondern Hinweise: Definition
    unscharf? Segment mehrdeutig?
  \ite Kategoriensystem überarbeiten, dokumentieren, neu codieren
\ei
\nsl{Für Abschlussarbeiten genügt meist die prozentuale Übereinstimmung
mit einer Diskussion der Abweichungen. Zufallskorrigierte Maße wie
Cohens $\kappa$ werden in Methodenlehrbüchern behandelt
\parencite{doeringbortz2016}.}
\os{\newpage}
\end{description}

%==============================================================================
\section{Stärken und Grenzen}
%==============================================================================
\begin{description}
\item[{\rot\bf Was qualitative Forschung leisten kann}]~\\[-\lblskip]
\bi
  \ite \textbf{stark}: Neues entdecken, Bedeutungen und Abläufe verstehen,
    Hypothesen bilden
  \ite \textbf{begrenzt}: Häufigkeiten schätzen, Ursachen prüfen,
    verallgemeinern
  \ite \textbf{Risiko}: Die Forschenden finden, was sie erwartet haben
    (LE~7)
  \ite \textbf{Gegenmittel}: offene Regeln, zweite Codierung, Triangulation
\ei
\nsl{Viele Projekte verbinden beide Zugänge. Ein typischer Weg führt von
Interviews, die Kategorien liefern, zu einem Fragebogen, der sie zählt.
Wie sich verschiedene Datenquellen gegenseitig absichern, ist Thema der
Triangulation in LE~10 \parencite{denzin1970research}.}
\os{\newpage}
\end{description}

\begin{merkbox}
Wer fragt, formt die Antwort mit. Qualitative Forschung leugnet das nicht,
sondern macht es sichtbar: durch offene Fragen, offengelegte Regeln und
eine zweite Person, die dasselbe Material unabhängig liest.
\end{merkbox}
\os{\newpage}

%==============================================================================
\section{Übung in der Stunde}
%==============================================================================
\begin{description}
\item[{\rot\bf Interview-Design (Einzeln und Paare, 25 Min.)}]~\\[-\lblskip]
\bi
  \ite Leitfaden zu Ihrer Forschungsfrage: 8--12 Fragen in der Vorlage
    (Aufgabe~2)
  \ite Probe-Interview in Paaren, 5 Minuten, dann Rollentausch
  \ite welche Frage war suggestiv, welche zu offen?
\ei
\os{\newpage}

\item[{\rot\bf Codier-Vergleich (zwei Gruppen, 15 Min.)}]~\\[-\lblskip]
\bi
  \ite beide Gruppen codieren dasselbe Transkript (Aufgabe~3)
  \ite jedes Segment genau einer Kategorie K1--K5 zuordnen
  \ite dann vergleichen: In wie vielen Segmenten stimmen Sie überein?
  \ite wo nicht -- und warum?
\ei
\os{\newpage}
\end{description}

\begin{lehrhinweis}
Beim Probe-Interview fallen fast immer dieselben Muster auf:
Ja/Nein-Fragen, Doppelfragen und Fragen, die eine Bewertung vorgeben.
Gute Reflexionen nennen ein konkretes Beispiel aus dem eigenen Leitfaden
und eine Umformulierung. Beim Codier-Vergleich ist das Transkript so
gebaut, dass einige Segmente eindeutig und einige absichtlich mehrdeutig
sind (Musterlösung im Erwartungshorizont des Übungsblatts). Volle
Übereinstimmung ist deshalb nicht zu erwarten. Daran lässt
sich zeigen, dass Intersubjektivität kein Zustand ist, sondern eine
Arbeit: Abweichungen besprechen, Kodierregeln schärfen, dokumentieren.
Lehrreich ist auch Segment~4: Dort stellt die interviewende Person eine
suggestive Frage, und die Antwort übernimmt deren Wortwahl.
\end{lehrhinweis}

%==============================================================================
\section{Lernziele und Aufgaben}
%==============================================================================
\begin{description}
\item[{\rot\bf Nach dieser Einheit können Sie}]~\\[-\lblskip]
\bi
  \ite qualitative und quantitative Forschungslogik voneinander abgrenzen
  \ite erklären, warum Sprache Daten mitformt
  \ite Leitfadeninterview, Fokusgruppe, Beobachtung und
    Dokumentenanalyse einordnen
  \ite einen offenen Interviewleitfaden entwerfen
  \ite die Gütekriterien nach Lincoln und Guba benennen und eine
    Inhaltsanalyse in Grundzügen durchführen
\ei
\os{\newpage}

\item[{\rot\bf Bis zur nächsten Woche}]~\\[-\lblskip]
\bi
  \ite \textbf{Interview-Design}: Leitfaden überarbeiten, ein
    Probe-Interview führen und kurz reflektieren (Aufgabe~2)
  \ite \textbf{Mini-Inhaltsanalyse}: Codier-Vergleich abschließen
    (Aufgabe~3)
  \ite \textbf{Tagebuch des Staunens}: Wann hat eine Formulierung Ihr
    Erleben verändert?
  \ite für LE~10: Notieren Sie, was Sie in Ihrem Projekt bisher
    \emph{gefunden} haben
\ei
\os{\newpage}
\end{description}

\handoutinput{handouts/LE09}
